% GENERATED FILE. Edit canonical lessons, then run npm run generate:books.
% canonical-source-hash: fnv1a64:276d86d85f329576
% canonical-lessons: KA-C417-kantu, KA-C417-baki, KA-C417-ulitaya, KA-C417-vasuli-madu, KA-C417-sallisu

\chapter{Paying Back}
\label{ch:ka-paying-back}
\hlchaptermodality{417}

\begin{chapteropening}
\textbf{By the end of this chapter:} \emph{I can say an instalment, arrears, savings, to recover dues and to submit.}

Complete the last lesson of chapter 417: I can say an instalment, arrears, savings, to recover dues and to submit.
\end{chapteropening}

\section[\texorpdfstring{\kn{ಕಂತು} (kantu)}{ಕಂತು (kantu)}]{\kn{ಕಂತು} (kantu) — an instalment}

\label{lesson:KA-C417-kantu}



\begin{quote}
Before the new one: say the Kannada for to ask formally, then the Kannada for an oath of office.
\end{quote}

\subsection*{You'll want to know: \kn{ಕಂತು}}
\textbf{\kn{ಕಂತು}} — \emph{kantu} — \textquotedblleft{}an instalment\textquotedblright{}.

\begin{quote}
\kn{ನಾನು} \kn{ಪ್ರತಿ} \kn{ತಿಂಗಳು} \kn{ಸಾಲದ} \kn{ಕಂತು} \kn{ಕಟ್ಟುತ್ತೇನೆ}.
\end{quote}

\begin{quote}
\emph{nānu prati tiṅgaḷu sālada kantu kaṭṭuttēne.}
\end{quote}

\begin{quote}
I pay an instalment of the loan every month.
\end{quote}

\kn{ಕಂತು} \kn{ಕಟ್ಟು}, to pay an instalment.

\begin{cousinweb}
\textbf{\kn{ಕಂತು}} (\emph{kantu}) is native Dravidian: DEDR 1209 lists Kannada \emph{kanta}, \emph{kantu}, \textquotedblleft{}an instalment, a fixed term of payment\textquotedblright{}.
\end{cousinweb}

\begin{grammarlens}[title={What you've built}]
One more word for saying what happens if.
\end{grammarlens}

\subsection*{Your turn}
\begin{itemize}
\raggedright
  \item \emph{Say it:} \emph{kantu}
  \item \emph{Say it:} \emph{kantu}, once more
  \item \emph{Say it:} \emph{nānu prati tiṅgaḷu sālada kantu kaṭṭuttēne.} — I pay an instalment of the loan every month.
  \item \emph{Recall:} say the Kannada for to ask formally, then the Kannada for an oath of office, then say \emph{kantu} again
\end{itemize}

\begin{tcolorbox}[breakable,colback=teal!4,colframe=teal!35!black,title={Before you move on}]
What does \kn{ಕಂತು} mean? (\textquotedblleft{}An instalment\textquotedblright{}.) Say it once more.
\end{tcolorbox}
\section[\texorpdfstring{\kn{ಬಾಕಿ} (bāki)}{ಬಾಕಿ (bāki)}]{\kn{ಬಾಕಿ} (bāki) — arrears; the balance due}

\label{lesson:KA-C417-baki}



\begin{quote}
Before the new one: say the Kannada for an oath of office, then the Kannada for an instalment.
\end{quote}

\subsection*{You'll want to know: \kn{ಬಾಕಿ}}
\textbf{\kn{ಬಾಕಿ}} — \emph{bāki} — \textquotedblleft{}arrears; the balance due\textquotedblright{}.

\begin{quote}
\kn{ಅಂಗಡಿಗೆ} \kn{ಇನ್ನೂ} \kn{ನೂರು} \kn{ರೂಪಾಯಿ} \kn{ಬಾಕಿ} \kn{ಇದೆ}.
\end{quote}

\begin{quote}
\emph{aṅgaḍige innū nūru rūpāyi bāki ide.}
\end{quote}

\begin{quote}
A hundred rupees are still owed to the shop.
\end{quote}

\kn{ಬಾಕಿ} \kn{ಇರು}, to be outstanding.

\begin{cousinweb}
\textbf{\kn{ಬಾಕಿ}} (\emph{bāki}) is a loan from Arabic \emph{bāqī}, \textquotedblleft{}remaining\textquotedblright{}, which reached Kannada through Hindustani.
\end{cousinweb}

\begin{grammarlens}[title={What you've built}]
One more word for saying what happens if.
\end{grammarlens}

\subsection*{Your turn}
\begin{itemize}
\raggedright
  \item \emph{Say it:} \emph{bāki}
  \item \emph{Say it:} \emph{bāki}, once more
  \item \emph{Say it:} \emph{aṅgaḍige innū nūru rūpāyi bāki ide.} — A hundred rupees are still owed to the shop.
  \item \emph{Recall:} say the Kannada for an oath of office, then the Kannada for an instalment, then say \emph{bāki} again
\end{itemize}

\begin{tcolorbox}[breakable,colback=teal!4,colframe=teal!35!black,title={Before you move on}]
What does \kn{ಬಾಕಿ} mean? (\textquotedblleft{}Arrears; the balance due\textquotedblright{}.) Say it once more.
\end{tcolorbox}
\section[\texorpdfstring{\kn{ಉಳಿತಾಯ} (uḷitāya)}{ಉಳಿತಾಯ (uḷitāya)}]{\kn{ಉಳಿತಾಯ} (uḷitāya) — savings}

\label{lesson:KA-C417-ulitaya}



\begin{quote}
Before the new one: say the Kannada for an instalment, then the Kannada for arrears.
\end{quote}

\subsection*{You'll want to know: \kn{ಉಳಿತಾಯ}}
\textbf{\kn{ಉಳಿತಾಯ}} — \emph{uḷitāya} — \textquotedblleft{}savings\textquotedblright{}.

\begin{quote}
\kn{ಉಳಿತಾಯದ} \kn{ಹಣದಿಂದ} \kn{ಅವಳು} \kn{ಸೈಕಲ್} \kn{ಕೊಂಡಳು}.
\end{quote}

\begin{quote}
\emph{uḷitāyada haṇadinda avaḷu saikal koṇḍaḷu.}
\end{quote}

\begin{quote}
She bought a bicycle with her savings.
\end{quote}

\kn{ಉಳಿತಾಯ} \kn{ಖಾತೆ}, a savings account.

\begin{cousinweb}
\textbf{\kn{ಉಳಿತಾಯ}} (\emph{uḷitāya}) is built on the Kannada verb \emph{uḷi}, \textquotedblleft{}to remain, to be left over\textquotedblright{}; DEDR does not list \emph{uḷitāya}, so the formation is not settled.
\end{cousinweb}

\begin{grammarlens}[title={What you've built}]
One more word for saying what happens if.
\end{grammarlens}

\subsection*{Your turn}
\begin{itemize}
\raggedright
  \item \emph{Say it:} \emph{uḷitāya}
  \item \emph{Say it:} \emph{uḷitāya}, once more
  \item \emph{Say it:} \emph{uḷitāyada haṇadinda avaḷu saikal koṇḍaḷu.} — She bought a bicycle with her savings.
  \item \emph{Recall:} say the Kannada for an instalment, then the Kannada for arrears, then say \emph{uḷitāya} again
\end{itemize}

\begin{tcolorbox}[breakable,colback=teal!4,colframe=teal!35!black,title={Before you move on}]
What does \kn{ಉಳಿತಾಯ} mean? (\textquotedblleft{}Savings\textquotedblright{}.) Say it once more.
\end{tcolorbox}
\section[\texorpdfstring{\kn{ವಸೂಲಿ} \kn{ಮಾಡು} (vasūli māḍu)}{ವಸೂಲಿ ಮಾಡು (vasūli māḍu)}]{\kn{ವಸೂಲಿ} \kn{ಮಾಡು} (vasūli māḍu) — to recover dues}

\label{lesson:KA-C417-vasuli-madu}



\begin{quote}
Before the new one: say the Kannada for arrears, then the Kannada for savings.
\end{quote}

\subsection*{You'll want to know: \kn{ವಸೂಲಿ} \kn{ಮಾಡು}}
\textbf{\kn{ವಸೂಲಿ} \kn{ಮಾಡು}} — \emph{vasūli māḍu} — \textquotedblleft{}to recover dues\textquotedblright{}.

\begin{quote}
\kn{ಬ್ಯಾಂಕ್} \kn{ಸಾಲವನ್ನು} \kn{ವಸೂಲಿ} \kn{ಮಾಡಿತು}.
\end{quote}

\begin{quote}
\emph{byāṅk sālavannu vasūli māḍitu.}
\end{quote}

\begin{quote}
The bank recovered the loan.
\end{quote}

Past (I …): \kn{ವಸೂಲಿ} \kn{ಮಾಡಿದೆ}, \emph{vasūli māḍide}.

\kn{ವಸೂಲಿ} alone is collection of money owed.

\begin{cousinweb}
\textbf{\kn{ವಸೂಲಿ} \kn{ಮಾಡು}} (\emph{vasūli māḍu}) is built from \emph{vasūli}, a loan from Hindustani \emph{vasūl}, \textquotedblleft{}collected\textquotedblright{}, which comes from Arabic \emph{wuṣūl}, \textquotedblleft{}arrival, receipt\textquotedblright{}, with \emph{māḍu}, \textquotedblleft{}to do\textquotedblright{}.
\end{cousinweb}

\begin{grammarlens}[title={What you've built}]
One more word for saying what happens if.
\end{grammarlens}

\subsection*{Your turn}
\begin{itemize}
\raggedright
  \item \emph{Say it:} \emph{vasūli māḍu}
  \item \emph{Say it:} \emph{vasūli māḍu}, once more
  \item \emph{Say it:} \emph{byāṅk sālavannu vasūli māḍitu.} — The bank recovered the loan.
  \item \emph{Recall:} say the Kannada for arrears, then the Kannada for savings, then say \emph{vasūli māḍu} again
\end{itemize}

\begin{tcolorbox}[breakable,colback=teal!4,colframe=teal!35!black,title={Before you move on}]
What does \kn{ವಸೂಲಿ} \kn{ಮಾಡು} mean? (\textquotedblleft{}To recover dues\textquotedblright{}.) Say it once more.
\end{tcolorbox}
\section[\texorpdfstring{\kn{ಸಲ್ಲಿಸು} (sallisu)}{ಸಲ್ಲಿಸು (sallisu)}]{\kn{ಸಲ್ಲಿಸು} (sallisu) — to submit; to pay}

\label{lesson:KA-C417-sallisu}



\begin{quote}
Before the new one: say the Kannada for savings, then the Kannada for to recover dues.
\end{quote}

\subsection*{You'll want to know: \kn{ಸಲ್ಲಿಸು}}
\textbf{\kn{ಸಲ್ಲಿಸು}} — \emph{sallisu} — \textquotedblleft{}to submit; to pay\textquotedblright{}.

\begin{quote}
\kn{ನಾನು} \kn{ಕಚೇರಿಗೆ} \kn{ಅರ್ಜಿ} \kn{ಸಲ್ಲಿಸಿದೆ}.
\end{quote}

\begin{quote}
\emph{nānu kacērige arji salliside.}
\end{quote}

\begin{quote}
I submitted an application to the office.
\end{quote}

Past (I …): \kn{ಸಲ್ಲಿಸಿದೆ}, \emph{salliside}.

Also to offer: \kn{ಧನ್ಯವಾದ} \kn{ಸಲ್ಲಿಸು}, to offer thanks.

\begin{cousinweb}
\textbf{\kn{ಸಲ್ಲಿಸು}} (\emph{sallisu}) is native Dravidian: DEDR 2781 lists Kannada \emph{sal}, \textquotedblleft{}to pass, to be current\textquotedblright{}, and its causative \emph{sallisu}.
\end{cousinweb}

\begin{grammarlens}[title={What you've built}]
One more word for saying what happens if.
\end{grammarlens}

\subsection*{Your turn}
\begin{itemize}
\raggedright
  \item \emph{Say it:} \emph{sallisu}
  \item \emph{Say it:} \emph{sallisu}, once more
  \item \emph{Say it:} \emph{nānu kacērige arji salliside.} — I submitted an application to the office.
  \item \emph{Recall:} say the Kannada for savings, then the Kannada for to recover dues, then say \emph{sallisu} again
\end{itemize}

\begin{tcolorbox}[breakable,colback=teal!4,colframe=teal!35!black,title={Before you move on}]
What does \kn{ಸಲ್ಲಿಸು} mean? (\textquotedblleft{}To submit; to pay\textquotedblright{}.) Say it once more.
\end{tcolorbox}
